\section{The Completed Determinant and Cyclotomic Twists}
\label{sec:completion-twists}

\subsection*{Base field and smooth projective models}

Throughout the geometric discussion we take \(k=\QQ\). Since
irreducibility, smoothness, and genus are unchanged after extending the
constant field to \(\overline{\QQ}\), we pass freely between \(k\) and
\(\overline{\QQ}\) when convenient. For an affine curve
\(C_{\mathrm{aff}}\subset\mathbb{A}^2_k\) defined by an irreducible
polynomial equation, we write \(\overline{X}\) (resp.\ \(\overline{Y}\))
for the smooth projective curve whose function field is
\(k(C_{\mathrm{aff}})\) (resp.\ its degree-two subfield under the
relevant involution). All genus and Hurwitz statements are made for
these complete nonsingular models.

\subsection{The unweighted spectrum}

\begin{proposition}[Explicit determinant]\label{prop:kernel-delta-explicit}
One has
\[
  \Delta(z,u)
  =1-(1+u)z-5uz^2+3u(1+u)z^3-u(u^2-3u+1)z^4
\]
\[
  \qquad
  +u(u^3-3u^2-3u+1)z^5+u^2(u^2+u+1)z^6.
\]
\end{proposition}

\begin{proof}
A symbolic determinant computation from the explicit \(10\times 10\)
matrix \(B(u)\) of \Cref{def:sync-kernel} gives the displayed formula;
the displayed polynomial is the exact certificate used in the sequel,
and the command block producing it is recorded in
\Cref{sec:appendix-computation}\,\textup{(a)}.
As a consistency check, the closed form satisfies the self-duality
identity of \Cref{prop:sync-self-duality}; in particular, comparing
\(\Delta(z,u)\) with \(\Delta(uz,u^{-1})\) pairs the coefficients of
\(z^k u^\ell\) and \(z^k u^{k-\ell}\) exactly as required.
\end{proof}

Although \(B(u)\) is a \(10\times 10\) matrix, the determinant has
degree \(6\) in \(z\) because four rows are redundant for every \(u\):
rows \(5,6,8,10\) are \(u\)-multiples of rows \(1,2,7,9\),
respectively. Hence \(\operatorname{rank} B(u)\le 6\), so \(0\) is a
persistent eigenvalue of multiplicity at least four.

\begin{proposition}[Untwisted Perron root and residue constant]
  \label{prop:sync-kernel-spectrum}
  The matrix \(B=B(1)\) is primitive and every row of \(B\) has sum
  \(3\). Consequently the spectral radius is determined before any
  characteristic-polynomial factorisation is used:
  \[
    B\mathbf 1=3\mathbf 1,\qquad
    \|B\|_\infty=\max_i\sum_j |B_{ij}|=3,\qquad
    \rho(B)\le \|B\|_\infty,
  \]
  so \(\rho(B)=3\), and the Perron eigenvalue \(3\) is simple and is
  the unique eigenvalue of modulus \(3\).
  In addition,
  \[
    \det(I-zB)=(z-1)(z+1)(3z-1)(z^3-z^2+z+1).
  \]
  The Perron residue constant \(C_B\) of
  \Cref{def:perron-residue-constant} is
  \[
    C_B=\frac{243}{272}.
  \]
\end{proposition}

\begin{proof}
Every state reaches \(\st{000}\) in at most three steps by the paths
listed in \Cref{def:sync-kernel}. Conversely, the transition table gives
explicit paths from \(\st{000}\) to the other states:
\[
  \st{000}\to\st{001},\qquad
  \st{000}\to\st{002},\qquad
  \st{000}\to\st{001}\to\st{010},
\]
\[
  \st{000}\to\st{001}\to\st{100},\qquad
  \st{000}\to\st{001}\to\st{101},\qquad
  \st{000}\to\st{002}\to\st{11-1},
\]
\[
  \st{000}\to\st{001}\to\st{010}\to\st{0-12},\qquad
  \st{000}\to\st{001}\to\st{010}\to\st{0-12}\to\st{01-1},
\]
\[
  \st{000}\to\st{002}\to\st{11-1}\to\st{1-12}.
\]
Thus the digraph of \(B\) is strongly connected. Since there is also the
self-loop \(\st{000}\to\st{000}\), the period is \(1\), so \(B\) is primitive.

We next determine the spectral radius directly from the non-negative
matrix \(B\), before using any determinant factorisation. Inspecting the
transition table at \(u=1\), every row of \(B=B(1)\) has sum \(3\).
Therefore
\[
  B\mathbf 1=3\mathbf 1,\qquad
  \|B\|_\infty=\max_i\sum_j |B_{ij}|=\max_i\sum_j B_{ij}=3,
  \qquad
  \rho(B)\le\|B\|_\infty.
\]
The first equality exhibits \(3\) as an eigenvalue. The inequality
follows from the \(\ell^\infty\)-operator norm, or equivalently from
\Cref{lem:row-sum-perron-bound}, and shows that no eigenvalue of \(B\)
can have modulus greater than \(3\). Hence \(3\le\rho(B)\le3\), so
\(\rho(B)=3\). Since \(B\) is primitive, the Perron--Frobenius theorem
gives that this eigenvalue is simple and that no other eigenvalue lies
on \(\abs{\lambda}=3\).

The determinant is used only after this Perron conclusion has been
proved, and only for the displayed factorisation and residue. Setting
\(u=1\) in
\Cref{prop:kernel-delta-explicit} gives
\[
  \det(I-zB)=(z-1)(z+1)(3z-1)(z^3-z^2+z+1).
\]
This factorisation is deliberately not used to prove the spectral
radius. In the variable \(z\), the cubic factor \(z^3-z^2+z+1\) has a
zero with \(\abs{z}<1\), so the false reciprocal-root argument would not
distinguish \(3z-1\) from all other factors. The Perron conclusion above
comes only from \(B\mathbf1=3\mathbf1\), the row-sum norm bound
\(\rho(B)\le\|B\|_\infty=3\), and primitivity. The factorisation is used
here solely for the residue computation
\[
  \begin{aligned}
  C_B
  &=\lim_{z\to1/3}
    \frac{1-3z}{(z-1)(z+1)(3z-1)(z^3-z^2+z+1)}  \\
  &=-\frac{1}{(-2/3)(4/3)(34/27)}
    =\frac{243}{272}.
  \end{aligned}
\]
This proves part~\textup{(i)} of
\Cref{thm:intro-kernel-geometry}.
\end{proof}


Together with \Cref{lem:primitive-orbit-asymptotic}, this yields
\[
  p_n(B)=\frac{3^n}{n}
  +O\!\left(\frac{\max\{\Lambda(B)^n,3^{n/2}\}}{n}\right).
\]


\subsection{Completion and algebraic geometry}

Introducing
\[
  u=r^2,\qquad s=r+r^{-1},\qquad w=zr,
\]
the self-duality of \(\Delta(z,u)\) implies that
\(\Delta(w/r,r^2)\) is a polynomial in \(w\) and \(s\), which we denote
by \(\widehat\Delta(w,s)\).

\begin{proposition}[Explicit completed determinant]
  \label{prop:sync-hatdelta-completion}
  The completed determinant is
  \[
    \widehat\Delta(w,s)
    =1-sw-5w^2+3sw^3+(5-s^2)w^4+(s^3-6s)w^5+(s^2-1)w^6.
  \]
  It satisfies
  \[
    \widehat\Delta(w,-s)=\widehat\Delta(-w,s).
  \]
\end{proposition}

\begin{proof}
Substitute \(u=r^2\) and \(z=w/r\) into the explicit formula of
\Cref{prop:kernel-delta-explicit}. Terms in \(r\) and \(r^{-1}\)
combine to give a polynomial in \(s=r+r^{-1}\), and a direct expansion
produces the displayed formula. This displayed formula, together with
the parity identity below it, is the certificate used in the geometric
and cyclotomic arguments; the command block producing it is
recorded in \Cref{sec:appendix-computation}\,\textup{(b)}. The parity
relation is immediate. This proves part~\textup{(ii)} of
\Cref{thm:intro-kernel-geometry}.
\end{proof}

\begin{remark}
The phrase ``completed determinant'' is non-standard. In this note it
simply denotes the polynomial obtained by rewriting \(\Delta(w/r,r^2)\)
in the self-dual invariants \(w\) and \(s=r+r^{-1}\).
\end{remark}



\begin{proposition}[Generic Galois group]
  \label{prop:sync-hatdelta-galois-s6-generic}
  Viewed as a polynomial in \(w\), the completed determinant
  \(\widehat\Delta(w,s)\) has generic Galois group \(S_6\) over
  \(\QQ(s)\).
\end{proposition}

\begin{proof}

Let \(G\le S_6\) be the generic Galois group of
\(\widehat\Delta(w,s)\) over \(\QQ(s)\). We identify \(G\) in four
steps.

\medskip
\noindent\textit{Step 1 (discriminant, \(G\not\subset A_6\)).}\quad
Write \(D(s):=\disc_w\widehat\Delta(w,s)\in\QQ[s]\). The explicit
SageMath computation recorded in
\Cref{sec:appendix-computation}\,(d) prints the output
\((20,-256)\) for \((\deg D,\operatorname{lc}D)\), so \(D(s)\) has
degree~\(20\) with leading coefficient~\(-256\).
This metadata is already a nonsquare certificate. Indeed, if a
polynomial \(D\in\QQ[s]\) is a square in \(\QQ(s)\), write
\(D=(A/B)^2\) with coprime \(A,B\in\QQ[s]\). Then
\(A^2=DB^2\), so every irreducible divisor of \(B\) divides \(A\);
hence \(B\) is constant and \(D\) is a square in \(\QQ[s]\). Its
leading coefficient is therefore a square in \(\QQ^\times\), in
particular positive. The leading coefficient \(-256\) rules this out.
Since the discriminant of \(\widehat\Delta(w,s)\) as a polynomial in
\(w\) changes sign under odd permutations,
\(D(s)\notin \QQ(s)^2\) implies \(G\not\subset A_6\).

\medskip
\noindent\textit{Step 2 (specialisation input).}\quad
We use the following standard fact: if \(s_0\in\QQ\) is such that
\(\widehat\Delta(w,s_0)\) is separable over \(\QQ\) (equivalently,
\(D(s_0)\neq 0\)), then the Galois group \(G_{s_0}\) of
\(\widehat\Delta(w,s_0)\) over \(\QQ\) embeds into \(G\); see the
specialisation theorem in \cite[\S 4.4, Proposition~1]{Serre2008Topics}.
The condition \(D(s_0)\neq 0\) is
verified at each chosen \(s_0\) by computing
\(\gcd(\widehat\Delta(w,s_0),\partial_w\widehat\Delta(w,s_0))\bmod p\)
for a convenient prime \(p\).

To read off cycle types inside \(G_{s_0}\), we apply Dedekind's
theorem in its localized form. Namely, let
\(f(w)\in\ZZ[w]\) have degree \(n\), and let \(p\) be a prime
which divides neither the leading coefficient of \(f\) nor
\(\disc(f)\). Then \(f\) has degree \(n\) after reduction modulo
\(p\), and multiplying its reduction by the inverse of the reduced
leading coefficient gives a monic polynomial with the same irreducible
factor degrees. If this reduction factors as
\(f_1\cdots f_r\) with the \(f_i\) distinct and irreducible over
\(\mathbb{F}_p\), then the Frobenius at \(p\) in \(\Gal(f/\QQ)\)
has cycle type \((\deg f_1,\dots,\deg f_r)\). This is the usual
Dedekind theorem \cite[I, \S{}6, Proposition~25]{Serre1979}, applied
over \(\ZZ_{(p)}\), in the Frobenius form used in
\cite[\S{}4.3]{Serre2008Topics}. The exact mod-\(p\) factorisation
data are recorded in \Cref{sec:appendix-computation}\,(d).

\medskip
\noindent\textit{Step 3 (three Frobenius specialisations).}

\smallskip
\noindent\textit{Transitivity (\(s=3\), mod \(7\)).}\quad
One checks that
\(\gcd(\widehat\Delta(w,3),\partial_w\widehat\Delta(w,3))\equiv 1\pmod 7\),
so \(D(3)\neq 0\) and the specialisation is separable.
Here \(\operatorname{lc}_w\widehat\Delta(w,3)=8\not\equiv0\pmod7\),
so the nonmonic normalization above is legitimate.
Modulo \(7\),
\[
  \widehat\Delta(w,3)
  \equiv w^6+2w^5+3w^4+2w^3+2w^2+4w+1\pmod 7,
\]
which is irreducible over \(\mathbb{F}_7\). Hence the Frobenius at \(7\)
in \(G_3\subset G\) is a \(6\)-cycle, so \(G\) acts transitively
on the six roots.

\smallskip
\noindent\textit{Transposition (\(s=2\), mod \(7\)).}\quad
One checks that
\(\gcd(\widehat\Delta(w,2),\partial_w\widehat\Delta(w,2))\equiv 1\pmod 7\),
so \(D(2)\neq 0\). Here
\(\operatorname{lc}_w\widehat\Delta(w,2)=3\not\equiv0\pmod7\).
Modulo \(7\),
\[
\begin{aligned}
  \widehat\Delta(w,2)
  &\equiv (w+6)(w+5)(w+2)(w+1) \\
  &\qquad\cdot(3w^2+3w+2)\pmod 7,
\end{aligned}
\]
The quadratic factor \(3w^2+3w+2\) is irreducible over
\(\mathbb{F}_7\). Thus the factorisation type is
\((1)(1)(1)(1)(2)\), so the Frobenius at \(7\) in \(G_2\subset G\) is a
transposition. Hence \(G\) contains a transposition.

\smallskip
\noindent\textit{5-cycle (\(s=3\), mod \(19\)).}\quad
One checks that
\(\gcd(\widehat\Delta(w,3),\partial_w\widehat\Delta(w,3))\equiv 1\pmod{19}\).
This squarefree reduction is exactly the separability witness:
equivalently \(19\nmid D(3)\), so Dedekind's theorem applies to this
factorisation.
Here \(\operatorname{lc}_w\widehat\Delta(w,3)=8\not\equiv0\pmod{19}\).
Modulo \(19\),
\[
  \widehat\Delta(w,3)
  \equiv (w+16)\bigl(8w^5+14w^4+9w^2+3w+6\bigr)\pmod{19},
\]
with the quintic factor irreducible over \(\mathbb{F}_{19}\). The
factorisation type is \((1)(5)\), so the Frobenius at \(19\) in
\(G_3\subset G\) is a \(5\)-cycle. Hence \(G\) contains a
\(5\)-cycle.

\medskip
\noindent\textit{Step 4 (subgroup classification).}\quad
We now have: \(G\le S_6\) is transitive, contains a transposition, and
contains a \(5\)-cycle. A transitive subgroup of \(S_6\) that contains
a \(5\)-cycle is primitive. Indeed, a non-trivial block system on six
points has either \(2\) or \(3\) blocks. Let \(\sigma\in G\) be a
\(5\)-cycle. The induced action of \(\sigma\) on such a block system
has order dividing both \(5\) and the order of a permutation group on
\(2\) or \(3\) points; hence the induced action is trivial. Therefore
each block is fixed setwise by \(\sigma\). Any block meeting the support
of the \(5\)-cycle must then contain the whole five-point support,
which is impossible for block size \(3\) or \(2\). Thus no non-trivial
block system exists, so \(G\) is primitive
\cite[Ch.~3, \S{}3.3]{Dixon1996}.
A primitive subgroup of \(S_6\) containing a transposition is the full
symmetric group \cite[Ch.~3, Theorem~3.3A]{Dixon1996}. Hence \(G=S_6\),
completing part~\textup{(iii)} of \Cref{thm:intro-kernel-geometry}.

\end{proof}

\begin{corollary}\label{lem:hatdelta-irreducible}
The polynomial \(\widehat\Delta(w,s)\) is irreducible in \(\QQ[w,s]\).
\end{corollary}

\begin{proof}
Since \(\Gal(\widehat\Delta/\QQ(s))=S_6\), the action on the six roots is
transitive, so \(\widehat\Delta(w,s)\) is irreducible over \(\QQ(s)\).
Gauss's lemma then implies irreducibility in \(\QQ[w,s]\).
\end{proof}
We first construct the quotient curve \(\overline{Y}\) and establish its
smoothness; then we record the integrality of \(C_{\mathrm{aff}}\) and
the non-squareness of \(x\) in \(k(\overline{Y})\), which together
justify the definition of \(\overline{X}\) as a well-defined smooth
projective curve.

\begin{proposition}[The quotient is a smooth plane quartic]
  \label{prop:sync-hatdelta-quotient}
  The invariants
  \[
    x:=w^2,\qquad y:=sw
  \]
  identify the sign quotient birationally with the projective closure of
  the affine curve
  \[
    F(x,y):=
    1-y-5x+3xy+5x^2-xy^2+xy^3-6x^2y+x^2y^2-x^3=0.
  \]
  This projective closure is smooth; consequently it is the smooth
  projective model \(\overline{Y}\), and
  \[
    g(\overline{Y})=3.
  \]
\end{proposition}

\begin{proof}

By \Cref{lem:self-dual-compression}, the coefficient of \(w^j\) in
\(\widehat\Delta(w,s)\) has degree at most \(j\) in \(s\) and the
parity of \(j\). Hence
\Cref{thm:sign-quotient-plane-genus} applies to the involution
\(\iota:(w,s)\mapsto(-w,-s)\): the invariants
\[
  x=w^2,\qquad y=sw
\]
identify the invariant function field birationally, and the theorem
gives a unique polynomial \(F(x,y)\in\QQ[x,y]\) with
\(\widehat\Delta(w,s)=F(w^2,sw)\).
A direct substitution gives the displayed formula for \(F\). Since
\(\widehat\Delta(0,s)=1\), the constant coefficient hypothesis
\(c_0\ne0\) in \Cref{thm:sign-quotient-plane-genus} is satisfied; hence
the quotient function field is the function field of \(F=0\). Thus
\(\overline{Y}\) is birational to the projective closure of the affine
quartic \(\{F=0\}\).
After the smoothness check below, this closure is the smooth projective
model \(\overline{Y}\).

We verify that this projective closure is smooth. Homogenizing
\(F(x,y)\) to degree \(4\) in projective coordinates \([X:Y:Z]\) gives
\[
\begin{aligned}
  F_h(X,Y,Z)
  &= Z^4 - YZ^3 - 5XZ^3 + 3XYZ^2 + 5X^2Z^2 - XY^2Z \\
  &\quad + XY^3 - 6X^2YZ + X^2Y^2 - X^3Z .
\end{aligned}
\]
On the affine chart \(Z=1\),
\[
  \partial_x F
  =-5+3y+10x-y^2+y^3-12xy+2xy^2-3x^2,
\]
\[
  \partial_y F
  =-1+3x-2xy+3xy^2-6x^2+2x^2y.
\]
A Gr\"obner-basis computation of the ideal
\(\langle F,\,\partial_x F,\,\partial_y F\rangle\subset\QQ[x,y]\)
gives \(\langle 1\rangle\), so there is no affine singular point.

The line at infinity \(Z=0\) meets the quartic where the leading form
\(XY^2(X+Y)=0\), giving the three points
\[
  [0:1:0],\qquad [1:0:0],\qquad [1:-1:0].
\]
At \([0:1:0]\) (chart \(Y=1\)): one computes
\[
\begin{aligned}
  \partial_X F_h\big|_{(0,1,0)}
  &=
  \bigl[-5Z^3+3YZ^2+10XZ^2-Y^2Z+Y^3 \\
  &\qquad -12XYZ+2XY^2-3X^2Z\bigr]_{(0,1,0)}
  = 1\neq 0 .
\end{aligned}
\]
At the other two points at infinity,
\[
  \partial_Z F_h\big|_{(1,0,0)}=-1\neq 0,
  \qquad
  \partial_Z F_h\big|_{(1,-1,0)}=4\neq 0 .
\]

Hence the projective curve is smooth everywhere, so
\Cref{thm:sign-quotient-plane-genus}, with \(n=4\), identifies this
nonsingular quartic as the smooth projective quotient model and gives
\[
  g(\overline{Y})=(4-1)(4-2)/2=3.
\]

\end{proof}

\begin{lemma}[Integrality and function-field setup]
  \label{lem:integral-function-field}
  Let \(k=\QQ\) and let \(\overline{Y}\) be the smooth projective quartic
  of \Cref{prop:sync-hatdelta-quotient}. Then:
  \begin{enumerate}[label=\textup{(\roman*)}]
    \item The affine curve \(C_{\mathrm{aff}}:=\{\widehat\Delta(w,s)=0\}
          \subset\mathbb{A}^2_k\) is integral.
    \item The element \(x=w^2\in k(\overline{Y})\) is not a square
          in \(k(\overline{Y})\).
    \item The extension \(k(\overline{Y})(\sqrt{x})/k(\overline{Y})\)
          has degree exactly two, and its function field is \(k(C_{\mathrm{aff}})\).
  \end{enumerate}
\end{lemma}

\begin{proof}

Part~(i): \(\widehat\Delta(w,s)\) is irreducible in \(\QQ[w,s]\) by
\Cref{lem:hatdelta-irreducible}, so its zero locus is an irreducible
variety. Since \(\widehat\Delta\) is prime in the UFD \(\QQ[w,s]\), the
coordinate ring \(\QQ[w,s]/(\widehat\Delta)\) is a domain, so the
variety is also reduced. Hence \(C_{\mathrm{aff}}\) is integral.

Part~(ii): From \Cref{prop:sync-hatdelta-quotient}, the affine quartic
satisfies \(F(0,y)=1-y\), so \(x=0\) meets \(\overline{Y}\) only at
the point \(P=(0,1)\), with \(\partial_y F(0,1)=-1\neq 0\), giving
\(\operatorname{ord}_P(x)=1\). A square in \(k(\overline{Y})\) would
have even valuation at every place; the odd order at \(P\) shows
\(x\notin k(\overline{Y})^2\).

Part~(iii): Since \(x\notin k(\overline{Y})^2\), the extension
\(k(\overline{Y})(\sqrt{x})/k(\overline{Y})\) has degree two. The
identifications \(w=\sqrt{x}\) and \(s=y/\sqrt{x}\) show that this
extension is the function field of \(C_{\mathrm{aff}}\).

\end{proof}

\begin{proposition}[Two-sheeted cover of genus six]
  \label{prop:sync-hatdelta-double-cover}
  The covering map \(\overline{X}\to\overline{Y}\) is a degree-two map
  branched at exactly two points of \(\overline{Y}\). Hence
  \[
    g(\overline{X})=6.
  \]
\end{proposition}

\begin{proof}

By \Cref{lem:integral-function-field}, the extension of function fields
\[
  k(\overline{X})/k(\overline{Y})
\]
is obtained by adjoining \(\sqrt{x}\), and it has degree exactly two.

We use the standard criterion for ramification in quadratic extensions
of function fields: let \(K\) be a function field of characteristic
\(\neq 2\), and let \(L=K(\sqrt{f}\,)\) with \(f\in K^\times\) not a
square. Then a place \(P\) of \(K\) ramifies in \(L/K\) if and only if
\(\operatorname{ord}_P(f)\) is odd; see \cite[III.7.3]{Stichtenoth2009}.
Here \(f=x\), and we determine the divisor of \(x\) on \(\overline{Y}\).

On the affine chart, \(F(0,y)=1-y\), so the only affine point with
\(x=0\) is \(P=(0,1)\). Since \(\partial_y F(0,1)=-1\neq 0\), the line
\(x=0\) meets \(\overline{Y}\) transversely there, and therefore
\[
  \operatorname{ord}_P(x)=1.
\]

For the points at infinity, work on the chart \(X=1\) with
\(y=Y/X\) and \(z=Z/X\). The equation becomes
\[
  G(y,z)=z^4-yz^3-5z^3+3yz^2+5z^2-y^2z+y^3-6yz+y^2-z.
\]
At \(Q_0=[1:0:0]\), one has \(\partial_z G(0,0)=-1\), so \(z\) is an
analytic function of \(y\), and solving \(G(y,z)=0\) gives
\[
  z=y^2-5y^3+O(y^4).
\]
Hence \(x=X/Z=1/z\) has pole order \(2\) at \(Q_0\), so
\[
  \operatorname{ord}_{Q_0}(x)=-2.
\]
At
\(Q_\infty=[1:-1:0]\), writing \(y=-1+t\) gives
\[
  G(-1+t,z)=t-2t^2+t^3+4z-4tz-t^2z
  +2z^2+3tz^2-4z^3-tz^3+z^4,
\]
so \(\partial_z G(-1,0)=4\neq 0\), and
\[
  z=-\frac14t+\frac{7}{32}t^2+O(t^3).
\]
Thus \(x=1/z\) has a simple pole at \(Q_\infty\), so
\[
  \operatorname{ord}_{Q_\infty}(x)=-1.
\]

At the remaining point \(Q_1=[0:1:0]\), use the chart \(Y=1\) with
\(u=X/Y\) and \(v=Z/Y\). The homogenised equation begins as
\[
  u-v^3+O(uv,v^4,u^2)=0,
\]
so \(u=v^3+O(v^4)\), and therefore \(x=X/Z=u/v\) has a zero of order
\(2\) at \(Q_1\), so
\[
  \operatorname{ord}_{Q_1}(x)=2.
\]

Thus
\[
  \operatorname{div}(x)=P+2Q_1-2Q_0-Q_\infty.
\]
The odd valuations occur exactly at \(P\) and \(Q_\infty\). By the
quadratic ramification criterion, these are exactly the two ramified
places of \(k(\overline{X})/k(\overline{Y})\), and each contributes one
simple branch point. The Riemann--Hurwitz genus formula
\cite[III.4.12]{Stichtenoth2009} then gives
\[
  2g(\overline{X})-2=2(2g(\overline{Y})-2)+2=2(2\cdot 3-2)+2,
\]
hence \(g(\overline{X})=6\).

\end{proof}

\begin{theorem}[Arithmetic ramification of the completed cover]
  \label{thm:sync-completed-ramification}
  Let
  \[
    \pi_s:\overline X\longrightarrow\mathbb P^1_s
  \]
  be the morphism induced by the completed parameter. Then \(\pi_s\)
  has degree \(6\). It has twenty finite branch values, with one simple
  ramification point over each, and two simple ramification points over
  \(s=\infty\). Thus the finite inertia type is
  \((2,1,1,1,1)\), the inertia type at infinity is
  \((2,2,1,1)\), and the ramification divisor has degree \(22\).

  The finite branch values are the roots of the irreducible polynomial
  \[
  \begin{aligned}
    D(s)={}&-256s^{20}+3712s^{18}-40320s^{16}+389241s^{14}
      -3214252s^{12}\\
    &+13200192s^{10}-20821228s^8+708704s^6
      +20467008s^4\\
    &-10514432s^2+147456.
  \end{aligned}
  \]
  Consequently they form a single
  \(\operatorname{Gal}(\overline{\mathbb Q}/\mathbb Q)\)-orbit. If
  \(D(s)=E(s^2)\), then
  \[
    \operatorname{Gal}(E/\mathbb Q)\cong S_{10};
  \]
  equivalently, the Galois action induced on the ten antipodal pairs of
  finite branch values is the full symmetric group.
\end{theorem}

\begin{proof}
By \Cref{lem:hatdelta-irreducible}, the function-field extension defined
by \(\widehat\Delta(w,s)\) has degree \(6\) over \(\mathbb Q(s)\), so
\(\deg\pi_s=6\). Direct elimination gives
\[
  \disc_w\widehat\Delta(w,s)=D(s)
\]
with \(D\) as displayed. This equality and all finite-field calculations
below are reproduced by the exact verifier
\[
  \texttt{python artifacts/verify\_ramification.py}.
\]

The reduction of \(D\) modulo \(71\) is
\[
\begin{aligned}
  28s^{20}+20s^{18}+8s^{16}+19s^{14}-11s^{12}+14s^{10}
  +19s^8\\
  {}-18s^6-20s^4+29s^2-11.
\end{aligned}
\]
It is irreducible in \(\mathbb F_{71}[s]\). Explicitly, modular
exponentiation gives
\[
  s^{71^{20}}\equiv s\pmod D,
\]
and
\[
  \gcd(s^{71^{10}}-s,D)=\gcd(s^{71^4}-s,D)=1
  \quad\text{in }\mathbb F_{71}[s].
\]
These are the irreducibility tests corresponding to the prime divisors
\(2\) and \(5\) of \(20\). Gauss's lemma therefore shows that \(D\) is
irreducible over \(\mathbb Q\), and hence squarefree. Moreover,
\(D(1)=D(-1)=325825\), so no root of \(D\) is a zero of the leading
coefficient \(s^2-1\) of \(\widehat\Delta\) in \(w\).

It follows from the discriminant criterion in characteristic zero that
each root of \(D\) has precisely one simple ramification point above it.
The only finite parameters at which a root can leave the affine
\(w\)-chart are \(s=\varepsilon\), where \(\varepsilon\in\{1,-1\}\).
With \(v=1/w\), put
\[
\begin{aligned}
  V(v,s):={}&v^6\widehat\Delta(v^{-1},s)\\
  ={}&v^6-sv^5-5v^4+3sv^3+(5-s^2)v^2
       +(s^3-6s)v+(s^2-1).
\end{aligned}
\]
At \((v,s)=(0,\varepsilon)\), one has
\(\partial_vV=-5\varepsilon\ne0\). Hence \(s-\varepsilon\) is a
uniformizer on the corresponding branch, and these two points are
unramified.

It remains to analyse the fibre over infinity. Set \(t=1/s\) and
\[
\begin{aligned}
 K(w,t):={}&t^3\widehat\Delta(w,t^{-1})\\
 ={}&t^3-wt^2-5w^2t^3+3w^3t^2+5w^4t^3-tw^4+w^5\\
    &\quad-6t^2w^5+tw^6-t^3w^6.
\end{aligned}
\]
The substitution \(w=ta\), followed by division by \(t^3\), has
special fibre \(1-a\). The implicit function theorem therefore gives
one unramified branch with \(w/t\to1\). For the remaining branches at
\((w,t)=(0,0)\), substitute \(t=vw^2\) and divide by \(w^5\). The
special fibre is
\[
  1-v^2.
\]
Its two simple roots \(v=1\) and \(v=-1\) give two branches
\[
  t=v_\pm(w)w^2,
  \qquad v_\pm(0)=\pm1,
\]
each with ramification index \(2\). Finally, for the branch with
\(w=\infty\), put \(q=1/w\). Then
\[
\begin{aligned}
 q^6K(q^{-1},t)
 ={}&t^3q^6-t^2q^5-5t^3q^4+3t^2q^3+5t^3q^2-tq^2\\
    &\quad+q-6t^2q+t-t^3,
\end{aligned}
\]
whose linear part at \((q,t)=(0,0)\) is \(q+t\). This is one further
unramified branch. The local degrees \(1,2,2,1\) sum to \(6\), so the
list is exhaustive. Thus infinity has inertia type \((2,2,1,1)\), and
the total ramification degree is \(20+2=22\).

For the final assertion, write
\[
\begin{aligned}
 E(x)={}&-256x^{10}+3712x^9-40320x^8+389241x^7-3214252x^6\\
 &+13200192x^5-20821228x^4+708704x^3+20467008x^2\\
 &-10514432x+147456.
\end{aligned}
\]
Modulo \(37\), this polynomial is irreducible, so its Galois group
contains a \(10\)-cycle. Modulo \(17\), it has squarefree factorisation
type \((1,9)\), so the group contains a \(9\)-cycle. Finally,
\[
  v_{1571}(\disc E)=1,
\]
and modulo \(1571\) its factorisation type, with multiplicities, is
\((1^2,2,6)\). The tame inertia group at \(1571\) therefore contains a
transposition.

The transitive group obtained from the irreducible reduction modulo
\(37\) is primitive. Indeed, a nontrivial block system on ten points
has blocks of size \(2\) or \(5\). A \(9\)-cycle cannot preserve two
blocks of size \(5\). If it permutes five blocks of size \(2\), then
its cube fixes every block and is a product of three \(3\)-cycles,
which is also impossible for blocks of size \(2\). A primitive subgroup
of \(S_{10}\) containing a transposition is \(S_{10}\)
\cite[Ch.~3, Theorem~3.3A]{Dixon1996}. Hence
\(\operatorname{Gal}(E/\mathbb Q)=S_{10}\), as claimed.
\end{proof}




